\documentclass[11pt]{article}
\usepackage[a4paper,margin=2.3cm,headheight=22pt]{geometry}
\usepackage[T1]{fontenc}
\usepackage{iftex}
\ifPDFTeX\usepackage[utf8]{inputenc}\fi
\usepackage{lmodern,microtype,booktabs,longtable,tabularx,array}
\usepackage{enumitem,xcolor,colortbl,titlesec}
\usepackage{xurl,hyperref,fancyhdr}
\definecolor{ink}{HTML}{183247}
\definecolor{accent}{HTML}{126C70}
\definecolor{muted}{HTML}{526676}
\definecolor{paperblue}{HTML}{EDF6F5}
\definecolor{gold}{HTML}{B68032}
\definecolor{papergold}{HTML}{FBF4E8}
\definecolor{rulecolor}{HTML}{B8D4D3}
\hypersetup{colorlinks=true,linkcolor=accent,urlcolor=accent}
\titleformat{\section}
  {\Large\bfseries\sffamily\color{ink}}
  {\colorbox{accent}{\color{white}\strut\thesection}}{0.6em}{}
  [\vspace{2pt}{\color{gold}\titlerule[0.7pt]}]
\titleformat{\subsection}
  {\large\bfseries\sffamily\color{accent}}{\thesubsection}{0.6em}{}
\titlespacing*{\section}{0pt}{12pt}{8pt}
\titlespacing*{\subsection}{0pt}{10pt}{5pt}
\arrayrulecolor{rulecolor}
\newcommand{\tablehead}{\rowcolor{paperblue}}
\setlength{\emergencystretch}{4em}
\setlength{\parindent}{0pt}
\setlength{\parskip}{6pt}
\setlist{itemsep=3pt,topsep=4pt}
\setcounter{tocdepth}{2}
\raggedbottom
\newcolumntype{L}[1]{>{\raggedright\arraybackslash}p{#1}}
\pagestyle{fancy}
\fancyhf{}
\lhead{\small\sffamily\textcolor{accent}{\textbf{CAMFRANGLAIS}}\;
  \textcolor{muted}{/ COMPILER CONSTRUCTION}}
\rhead{\small\sffamily\textcolor{muted}{CS4110 / SET A}}
\cfoot{\colorbox{paperblue}{\makebox[1.8em]{\small\sffamily\textcolor{accent}{\thepage}}}}
\renewcommand{\headrulewidth}{0.6pt}
\renewcommand{\headrule}{\color{accent}\hrule width\headwidth height\headrulewidth\vskip-\headrulewidth}
\newcommand{\notice}[2]{%
  \par\medskip\noindent\begingroup\setlength{\fboxsep}{8pt}%
  \colorbox{paperblue}{%
    \begin{minipage}{\dimexpr\linewidth-2\fboxsep\relax}
      {\bfseries\sffamily\textcolor{accent}{#1}}\par
      #2
    \end{minipage}}\endgroup\par\medskip}
\newcommand{\projectversion}{4.1}
\newcommand{\projectdate}{27 September 2026}
\newcommand{\projectcoverlabel}{DETERMINISTIC / PUBLIC ACCESS / EVIDENCE-LED}
\newcommand{\projectcover}[3]{%
  \hypersetup{pageanchor=false}
  \begin{titlepage}
    \noindent\colorbox{ink}{%
      \begin{minipage}{\dimexpr\linewidth-2\fboxsep\relax}
        \vspace{8pt}\hspace{8pt}%
        {\small\bfseries\sffamily\color{white}CS4110 \quad / \quad SUMMER 2026 \quad / \quad SET A}
        \par\vspace{8pt}
      \end{minipage}}\par
    \vspace{1.0cm}
    {\color{gold}\rule{3.4cm}{4pt}}\par
    \vspace{0.6cm}
    {\Huge\bfseries\color{ink}#1\par}
    \vspace{0.9cm}
    {\LARGE\sffamily\color{accent}Camfranglais\par}
    \vspace{0.4cm}
    {\large\color{muted}#2\par}
    \vspace{1cm}
    \notice{\projectcoverlabel}{#3}
    \vfill
    {\small\bfseries\sffamily\color{accent}AUTHOR GROUP}\par
    {\renewcommand{\arraystretch}{1.45}
    \begin{tabularx}{\linewidth}{@{}X r@{}}
      \toprule\tablehead\textbf{Name} & \textbf{Matricule}\\\midrule
      Kwete Ngouba Junior Rayan & ICTU20241377\\
      Djemtchimo Noukui Bruno Jonatan & ICTU20241585\\
      Amina Boubakary & ICTU20241870\\
      \bottomrule
    \end{tabularx}}
    \par\vspace{0.7cm}
    {\small\textcolor{accent}{\textbf{Version \projectversion}} \quad|\quad \projectdate\par
    Faculty of Information and Communication Technologies\par
    Instructor: Engr. Tanwi Nkiamboh\par
    Assignment due: 29 September 2026}
  \end{titlepage}
  \hypersetup{pageanchor=true}}

\renewcommand{\projectversion}{4.3}
\renewcommand{\projectdate}{28 September 2026}
\renewcommand{\projectcoverlabel}{STUDY OVERVIEW}
\usepackage{graphicx}
\newcommand{\CorpusTokens}{81}
\newcommand{\CorpusForms}{55}
\newcommand{\CorpusAccepted}{10}
\newcommand{\CorpusRejected}{2}
\newcommand{\CorpusSwitches}{26}
\newcommand{\DictionaryCount}{938}
\newcommand{\ReferenceRows}{878}
\newcommand{\GrammarProductions}{39}
\newcommand{\GrammarCells}{47}

\hypersetup{pdftitle={Camfranglais Final Coursework Report},
  pdfauthor={Kwete Ngouba Junior Rayan; Djemtchimo Noukui Bruno Jonatan; Amina Boubakary}}
\setcounter{tocdepth}{1}
\makeatletter
\renewcommand{\l@section}{\@dottedtocline{1}{0em}{2.5em}}
\makeatother
\newcommand{\code}[1]{\texttt{#1}}
\newcommand{\screen}[3]{%
  \par\smallskip
  {\centering
    \includegraphics[width=\linewidth,height=#2,keepaspectratio]{evidence/screenshots/#1}
    \par}
  {\small\textbf{\textcolor{accent}{Screenshot.}} #3\par}}
\newcommand{\automaton}[3]{%
  \par\smallskip
  {\centering
    \includegraphics[width=\linewidth,height=#2,keepaspectratio]{diagrams/#1.png}
    \par}
  \refstepcounter{figure}\label{fig:#1}%
  {\small\textbf{\textcolor{accent}{Figure \thefigure.}} #3\par}}
\begin{document}
\projectcover{Final Coursework Report}
{Lexical and Syntactic Analysis of Informal Urban Communication in Yaound\'e}
{We study twelve unchanged transcriptions with our regex lexer and
corpus-derived LL(1) parser. We present the original wording, lexical results,
slang and topic readings, grammar calculations, automata, complete traces and
screenshots of the working application.}
\setcounter{page}{2}

\section*{Abstract and Evidence Statement}
\addcontentsline{toc}{section}{Abstract and Evidence Statement}
We built Camfranglais to make the lexical and syntactic analysis of informal
expressions visible. We preserve spelling and word order, split each statement
with regular expressions, assign lexical categories and test the resulting
category sequence with an LL(1) parser.

We evaluate the twelve transcriptions S01--S12 exactly as stored in Collection.
One controlled pass produces \CorpusTokens{} tokens and \CorpusForms{}
case-folded forms. Our original grammar has ten nonterminals and 37
alternatives. Removing direct left recursion and applying left factoring
produce twelve nonterminals, \GrammarProductions{} alternatives and
\GrammarCells{} populated predictive-table cells, with no conflict.
Ten statements parse completely. S09 and S10 reject after five consumed
tokens, at the unchanged forms \code{n'ais} and \code{alli}.

We also discuss the difference between an emitted software label and a word's
role in its sentence. Our linguistic reading identifies 15 slang forms in
17 occurrences, while the unchanged lexer emits zero SLANG labels for this
corpus. We relate each statement to its main topic using the words and
French gloss recorded with it.

The application opens without an account. Visitors can analyze and save tests,
then inspect their original results. We use genuine browser captures from
27 September 2026 to illustrate this workflow.

\begin{tabularx}{\linewidth}{@{}X r@{}}
\toprule\tablehead\textbf{Controlled evaluation} & \textbf{Measured result}\\\midrule
Raw statements / original token occurrences & 12 / \CorpusTokens\\
Known-category occurrences / unresolved occurrences & 79 / 2\\
Emitted language-transition clues & \CorpusSwitches\\
Accepted / rejected category sequences & \CorpusAccepted{} / \CorpusRejected\\
LL(1) conflicts / residual left-recursive nonterminals & 0 / 0\\
Additional constructed boundary controls & 7\\
\bottomrule
\end{tabularx}

\notice{Reading the results}{
We use the 83.3\% grammar-match rate to describe these twelve category
sequences; it does not measure language-wide accuracy.}

\clearpage
{\setlength{\parskip}{0pt}\tableofcontents}

\clearpage
\section{Context, Problem and Objectives}
\subsection*{The assignment and the problem}
CS4110 SET A asks a group of three to examine 10--15 real-life expressions
heard in Yaound\'e, preserving slang, code-mixing, incomplete speech and
nonstandard forms. The required compiler work includes token categories and
regular expressions, frequency and variation analysis, a grammar constructed
from the expressions, grammar transformations, FIRST/FOLLOW, a parsing table
and accepted/rejected test cases. Our chosen parsing route is LL(1);
an LR automaton is not additionally required for that route.

We face a different problem from compiling a conventional programming
language: everyday speech depends on context and varies between speakers.
In our corpus, \code{school} occurs within a French-framed expression.
We therefore preserve the wording first and discuss spelling and meaning
separately from the computed categories.

Ebongue's study distinguishes speaker-written Camfranglais from researcher
transcriptions and discusses variation in spelling, syntax and morphosyntax
\cite{ebongue}. Stored text and computed categories are observable, but a
small textual sample cannot establish how a community speaks.

\subsection*{Operational objectives}
\begin{enumerate}
\item Preserve all twelve original strings and label every token in source order.
\item Apply deterministic recognition rules and expose unknowns, ambiguous
labels, language clues and corpus frequencies.
\item Derive reusable category productions from observed structures rather
than hard-code a whitelist of twelve complete sentences.
\item Perform grammar transformations and compute every set/table cell.
\item Demonstrate predictive parsing, successful consumption and precise
failure positions, including examples outside our twelve-statement corpus.
\item Present an accessible public workspace with durable, immutable test
evidence and no unnecessary account workflow.
\end{enumerate}

\subsection*{Scope}
We concentrate on lexical analysis, category-level syntax and visible
intermediate results. Translation, speech recognition and semantic parsing
are outside this study.

\clearpage
\section{Corpus Entry, Preparation and Ethics}
\subsection*{Entering the statements}
We first entered our twelve-statement corpus into the compiler for testing on
\textbf{22 September 2026}. We credit this corpus-entry work jointly to
\textbf{Kwete Ngouba Junior Rayan},
\textbf{Djemtchimo Noukui Bruno Jonatan} and
\textbf{Amina Boubakary}. The entry date refers to our compiler testing
session.

We use identifiers S01--S12 throughout the report. Each record keeps the
original statement beside its French gloss, so we can compare the wording,
its interpretation and the compiler's output.

\begin{tabularx}{\linewidth}{@{}L{4.4cm}X@{}}
\toprule\tablehead\textbf{Preparation stage} & \textbf{Our record}\\\midrule
Initial compiler input & 22 September 2026; twelve statements processed
in the initial local analysis.\\
Report analysis snapshot & 27 September 2026; the twelve report statements
and the grammar used for our final calculations.\\
Raw text & Capitalization, apostrophes, spelling and word order preserved.\\
Meanings and topics & French glosses kept separately; topic readings
explained through the exact words in each statement.\\
\bottomrule
\end{tabularx}

The initial test entry used \code{ali} in S10; our report keeps \code{alli},
as recorded in the 27 September snapshot. We use that snapshot consistently
for every token table and parser trace.

\subsection*{Working with informal spelling}
We had to separate the written form from our interpretation of it.
For example, we kept \code{sharp} and \code{shap} as distinct spellings
and preserved the contractions in \code{J'ai} and \code{n'est}.
We also left \code{n'ais} and \code{alli} unchanged so that their effect
on lexical recognition and parsing remains visible.

\subsection*{Preserving the corpus}
We checked every raw statement against its regression fixture. Our
27 September analysis snapshot contains the same twelve Sentence entries
and one separate Word annotation, \code{choa}. We exclude that Word entry
from the sentence and token totals.

For each statement, we retain all tokens in source order, calculate
frequencies and inspect the parser trace. Section~23 brings these formal
results back to everyday themes such as rest, school, bargaining,
electricity and transport.

\subsection*{Ethical handling}
We refer to statements by their identifiers and omit speaker names.
We keep raw wording, meanings and analytical notes separate. The public
application warns that submitted test text is visible and retained, so we
exclude personal or confidential information from demonstrations.

\clearpage
\section{Raw Corpus: Statements S01--S06}
We use report-local identifiers for the statements below. Case, apostrophes
and wording are exactly those saved in Collection. The French glosses help
us explain our reading of the expressions.

{\small\renewcommand{\arraystretch}{1.6}\par\noindent
\begin{tabularx}{\linewidth}{@{}L{0.7cm} >{\raggedright\arraybackslash}X >{\raggedright\arraybackslash}X@{}}
\toprule
\textbf{ID} & \textbf{Exact raw statement} & \textbf{Supplied French meaning}\\
\midrule
S01 & \texttt{J'ai n'est pas tchop depuis bayar} & Je n'ai pas mangé depuis longtemps.\\
S02 & \texttt{je suis kass je go nang} & Je suis fatigué, je vais dormir.\\
S03 & \texttt{je suis back du school le sharp man} & Je suis revenu de l'école ce matin.\\
S04 & \texttt{j'ai mban gars je yabat} & J'ai échoué, mon gars, j'ai mal.\\
S05 & \texttt{elles yamo flop le kongossa} & Elles perdent leur temps à commérer.\\
S06 & \texttt{go au piol shoua mes kako} & Va à la maison laver mes habits.\\
\bottomrule
\end{tabularx}\par
}

\subsection*{Reading the evidence}
We read \code{tchop} in S01 as an eating verb because the gloss contains
``mang\'e''. The current program emits NOUN, so we identify this as a
questionable lexical label while keeping the computed table unchanged.

In S02, we read fatigue followed by an intention to sleep. The two clauses
are placed next to each other. S03 combines a French frame with
\code{back}, \code{school}, \code{sharp} and \code{man}. S04 expresses a
setback and distress in its gloss. S05 concerns gossip, and S06 gives two
successive household instructions.

These readings help us relate the category grammar to communication.
The glosses can paraphrase the full expression rather than align with
every word individually.

\clearpage
\section{Raw Corpus: Statements S07--S12}
We complete our fixed corpus with the six records below. We calculate the
report totals from S01--S12 only, keeping public test history, dictionary
entries and practice examples in their separate datasets.

{\small\renewcommand{\arraystretch}{1.6}\par\noindent
\begin{tabularx}{\linewidth}{@{}L{0.7cm} >{\raggedright\arraybackslash}X >{\raggedright\arraybackslash}X@{}}
\toprule\tablehead
\textbf{ID} & \textbf{Exact raw statement} & \textbf{French gloss}\\
\midrule
S07 & \texttt{pere ca me chop je suis kass} & Père, ça me touche, je suis épuisé.\\
S08 & \texttt{Ont a cote la light depuis le shap} & On a coupé le courant depuis ce matin.\\
S09 & \texttt{shiba le price pere je n'ais pas flop les do} & Baisse le prix, s'il te plaît, je n'ai pas beaucoup d'argent.\\
S10 & \texttt{mon combi came au school alli day a pied} & Mon ami est venu à l'école, il est venu à pied.\\
S11 & \texttt{je falla le taco depuis bayar} & Je cherche le taxi depuis longtemps.\\
S12 & \texttt{les talk free le francais} & Ils parlent le français librement.\\
\bottomrule
\end{tabularx}\par
}

\subsection*{Interpretive themes}
We read fatigue or personal state in S07, electricity interruption in S08,
bargaining in S09, school and walking in S10, transport in S11, and language
use in S12. Section~23 records the exact cues behind these interpretations.

We preserve \code{Ont a cote} in S08. In context, we read \code{a} as a
French auxiliary, supported by the gloss ``On a coup\'e le courant''.
The unchanged lexer emits EF; our discussion makes that mismatch explicit.
S09 retains \code{n'ais}, and S10 retains \code{alli}. Both remain unresolved
and cause the first grammar failure in their respective statements.

\subsection*{Sampling boundary}
Our twelve records meet the assignment's numerical range. We use this small
development corpus both to construct and to exercise the grammar; an
independent sample would be needed to assess wider coverage.

\clearpage
\section{Custom Lexical Specification}
\subsection*{Segmentation by regular expression}
We use a tokenizer that recognizes Unicode word-like runs, internal apostrophes/hyphens,
decimal numbers, selected punctuation and finally any remaining non-whitespace
character. The final alternative is important: unsupported input is preserved
for classification rather than silently discarded.

The following factored notation describes the segmentation pattern.
\code{M} denotes the combining-mark ranges U+0300--036F, U+1AB0--1AFF,
U+1DC0--1DFF, U+20D0--20FF and U+FE20--FE2F. Substituting those ranges and
\code{L} into the definitions yields the pattern in the source; spaces around
the displayed alternatives are notation, not characters to match.
{\small
\begin{verbatim}
L     = [^\W\d_]+(?:[M]+[^\W\d_]*)*
WORD  = L(?:['\u2018\u2019\u02bc-]L)*
TOKEN = WORD | \d+(?:[.,]\d+)? | [.,!?;:"()] | \S
\end{verbatim}}

Internal apostrophes therefore keep \code{J'ai}, \code{n'est} and
\code{n'ais} as single tokens. Our twelve statements have no explicit
punctuation tokens, but punctuation behavior is exercised by separate controls.
Whitespace separates matches; the original full input, including its whitespace,
is preserved independently of the token list.

\subsection*{Regex classification and fallback}
{\small
\begin{verbatim}
NUMBER       ^\d+(?:[.,]\d+)?$
PUNCTUATION  ^[.,!?;:"()]+$
ENGLISH_VERB_LIKE  suffix "ing" or "ed"
FRENCH_VERB_LIKE   suffix "er", "ir" or "re", length > 3
\end{verbatim}}

We apply the first two rules as full matches before lookup. We use the suffix
checks on normalized words only after the explicit lexical sources have
failed. We retain the \code{\_LIKE} suffix to make a morphological guess
visible; a verbal-looking ending can have a different function in context.

\subsection*{Deterministic pipeline}
We tokenize the raw input, classify each token by the first applicable rule,
retain its spelling and category, add phrase and language clues, compute
statistics, and pass the categories to the parser. We keep the same input
throughout these stages. The implementation is in
\nolinkurl{compiler/lexer/tokenizer.py} and
\nolinkurl{compiler/lexer/lexicon.py} \cite{source}.

\clearpage
\section{Finite-State Model of the Lexer}
We model the regular token pattern with an equivalent deterministic finite
automaton (DFA). Figure~\ref{fig:automaton-lexer} explains the pattern executed
by Python's \code{re} matcher. A double circle accepts a \emph{token form};
vocabulary classification and full-statement parsing follow later.

\automaton{automaton-lexer}{8.2cm}
{Equivalent token recognizer, initially in \(q0\). Double-circle states accept:
\(W\) word; \(I\) integer; \(F\) fraction; \(P\) punctuation; \(O\) other character.
\(J\) awaits a joiner continuation and \(R\) a decimal digit; neither accepts.
Unshown transitions enter an implicit dead state.}

\textbf{Input classes.} \(C\) is \code{[\textasciicircum{}\string\W\string\d\_]}:
a Unicode word character excluding decimal digits and underscore.
\(D\) is \code{\string\d}; \(M\) is the combining-mark set on the preceding
page. Pending joiners are ASCII/curly apostrophes or hyphen, outside \(C\).
U+02BC is already a word character and follows the \(C\) edges, even at a word
boundary. Punctuation is \code{[.,!?;:"()]}; ``other'' is the remaining
non-whitespace input, including a standalone underscore or hyphen.

\textbf{Token boundaries.} For this pattern, greedy matching retains the
longest accepted prefix. A dangling joiner or decimal separator cannot finish
in \(J\) or \(R\), so the match ends at the previous accepting state and the
remaining character is scanned again. Whitespace is skipped between matches,
not emitted as a token. The source's final \code{\string\S} alternative
preserves other characters instead of losing them.

{\small\par\noindent
\begin{tabularx}{\linewidth}{@{}L{3.0cm} >{\raggedright\arraybackslash}X@{}}
\toprule\tablehead
\textbf{Raw input} & \textbf{Observed tokenization}\\
\midrule
\texttt{j'ai} & \texttt{["j'ai"]}\\
\texttt{go-slow} & \texttt{["go-slow"]}\\
\texttt{a-} & \texttt{["a", "-"]}\\
\texttt{3.14} & \texttt{["3.14"]}\\
\texttt{3.} & \texttt{["3", "."]}\\
\texttt{.5} & \texttt{[".", "5"]}\\
\texttt{-5} & \texttt{["-", "5"]}\\
\texttt{\_x} & \texttt{["\_", "x"]}\\
\bottomrule
\end{tabularx}\par
}

We executed these eight constructed checks through the real tokenizer.
For example, \code{3.} becomes a number and punctuation, while \code{-5}
becomes a hyphen and a number. A recognized word form can still receive
\code{UNKNOWN} during classification.

\clearpage
\section{Lexicon, Category Policy and Notation}
\subsection*{Precedence is part of the result}
We apply NUMBER/PUNCTUATION first, followed by compatible approved Word
annotations. Conflicting reviewed labels are left unresolved. Our approved
\code{choa} annotation occurs outside these twelve sentences.

Explicit lookup then proceeds through SLANG, PIDGIN\_MARKER, NOUN, VERB,
French function words and English function words. The reference CSV fills
remaining gaps; one category is retained, multiple candidates yield
AMBIGUOUS. Explicit French forms supplement gaps, then morphological guesses
and finally UNKNOWN apply. We keep this precedence visible when comparing a software label with our
reading of a word in context.

Our reference dictionary contains \DictionaryCount{} display entries from
\ReferenceRows{} CSV rows and 60 additional senses. We use normalized aliases
for lookup while keeping displayed spelling. Phrase annotations remain
separate from the individual token categories.

\begin{tabularx}{\linewidth}{@{}l X l X@{}}
\toprule\tablehead\textbf{Code} & \textbf{Lexer category} &
\textbf{Code} & \textbf{Lexer category}\\\midrule
FF & FRENCH\_FUNCTION\_WORD & EF & ENGLISH\_FUNCTION\_WORD\\
N & NOUN & V & VERB\\
ADJ & ADJECTIVE & ADV & ADVERB\\
AMB & AMBIGUOUS & PID & PIDGIN\_MARKER\\
PUNC & PUNCTUATION & UNK & UNKNOWN\\
SL & SLANG & NUM & NUMBER\\
\bottomrule
\end{tabularx}

We abbreviate labels only in the printed tables. The lexer supports twenty
categories; nine occur here. The grammar also uses nine terminals, replacing
the corpus's UNK with optional PUNC.

\subsection*{Slang register and grammatical role}
We identify the following 15 forms, occurring 17 times, as slang-register
readings in our corpus. A slang word can also serve as a verb, noun or
adjective. The application is unchanged: its emitted SLANG count remains
zero, and all token tables below retain their actual labels.

{\small
\begin{tabularx}{\linewidth}{@{}L{4.4cm}l X@{}}
\toprule\tablehead\textbf{Form(s)} & \textbf{Output} & \textbf{Our reading in context}\\\midrule
\code{nang, mban, yabat, yamo, shoua, falla} & V &
Verbal roles: sleeping, a setback, distress, attitude, washing and seeking.\\
\code{piol, kako, taco, combi} & N &
Nominal roles: home, clothes, taxi and friend in the glosses.\\
\code{kass} & ADJ & Fatigue in S02 and S07 supports an adjective reading.\\
\code{bayar} & ADV & With \code{depuis}, it expresses a time span.\\
\code{tchop} & N & S01's eating meaning supports a verb reading.\\
\code{kongossa} & AMB & After \code{le} in S05, gossip is a nominal complement.\\
\code{do} & EF & S09's \code{les do} refers to money, a noun reading.\\
\bottomrule
\end{tabularx}}

\clearpage
\section{Token Tables: S01--S04}
We show every emitted token in its original order. We use the preceding
legend and preserve capitalization and apostrophes in the raw-token column.

{\small\begin{minipage}[t]{0.48\linewidth}
\textbf{S01}\hfill 6 tokens

{\raggedright\texttt{J'ai n'est pas tchop depuis bayar}\par}

\par\noindent
\begin{tabularx}{\linewidth}{@{}r X l@{}}
\toprule\tablehead
\textbf{No.} & \textbf{Raw token} & \textbf{Label}\\
\midrule
1 & \texttt{J'ai} & \texttt{V}\\
2 & \texttt{n'est} & \texttt{V}\\
3 & \texttt{pas} & \texttt{FF}\\
4 & \texttt{tchop} & \texttt{N}\\
5 & \texttt{depuis} & \texttt{FF}\\
6 & \texttt{bayar} & \texttt{ADV}\\
\bottomrule
\end{tabularx}\par

\end{minipage}\hfill
\begin{minipage}[t]{0.48\linewidth}
\textbf{S02}\hfill 6 tokens

{\raggedright\texttt{je suis kass je go nang}\par}

\par\noindent
\begin{tabularx}{\linewidth}{@{}r X l@{}}
\toprule\tablehead
\textbf{No.} & \textbf{Raw token} & \textbf{Label}\\
\midrule
1 & \texttt{je} & \texttt{FF}\\
2 & \texttt{suis} & \texttt{V}\\
3 & \texttt{kass} & \texttt{ADJ}\\
4 & \texttt{je} & \texttt{FF}\\
5 & \texttt{go} & \texttt{V}\\
6 & \texttt{nang} & \texttt{V}\\
\bottomrule
\end{tabularx}\par

\end{minipage}\par\bigskip
\begin{minipage}[t]{0.48\linewidth}
\textbf{S03}\hfill 8 tokens

{\raggedright\texttt{je suis back du school le sharp man}\par}

\par\noindent
\begin{tabularx}{\linewidth}{@{}r X l@{}}
\toprule\tablehead
\textbf{No.} & \textbf{Raw token} & \textbf{Label}\\
\midrule
1 & \texttt{je} & \texttt{FF}\\
2 & \texttt{suis} & \texttt{V}\\
3 & \texttt{back} & \texttt{AMB}\\
4 & \texttt{du} & \texttt{FF}\\
5 & \texttt{school} & \texttt{N}\\
6 & \texttt{le} & \texttt{FF}\\
7 & \texttt{sharp} & \texttt{ADV}\\
8 & \texttt{man} & \texttt{AMB}\\
\bottomrule
\end{tabularx}\par

\end{minipage}\hfill
\begin{minipage}[t]{0.48\linewidth}
\textbf{S04}\hfill 5 tokens

{\raggedright\texttt{j'ai mban gars je yabat}\par}

\par\noindent
\begin{tabularx}{\linewidth}{@{}r X l@{}}
\toprule\tablehead
\textbf{No.} & \textbf{Raw token} & \textbf{Label}\\
\midrule
1 & \texttt{j'ai} & \texttt{V}\\
2 & \texttt{mban} & \texttt{V}\\
3 & \texttt{gars} & \texttt{PID}\\
4 & \texttt{je} & \texttt{FF}\\
5 & \texttt{yabat} & \texttt{V}\\
\bottomrule
\end{tabularx}\par

\end{minipage}\par\bigskip
}

\subsection*{What these classifications show}
We treat S01's \code{J'ai} and \code{n'est} as whole VERB tokens. Our grammar
therefore works with coarser units than a detailed analysis that separates
pronouns and auxiliaries.

We obtain two verbs from S02's \code{go nang}, enabling a serial-verb tail.
S03 contains two ambiguous tokens without an UNKNOWN. In S04, we obtain
PID for \code{gars} because the explicit Pidgin list has priority.

The lengths are six, six, eight and five tokens respectively. These 25
occurrences contribute to the corpus total once each.

\clearpage
\section{Token Tables: S05--S08}
We apply the same classification order to these four statements.

{\small\begin{minipage}[t]{0.48\linewidth}
\textbf{S05}\hfill 5 tokens

{\raggedright\texttt{elles yamo flop le kongossa}\par}

\par\noindent
\begin{tabularx}{\linewidth}{@{}r X l@{}}
\toprule\tablehead
\textbf{No.} & \textbf{Raw token} & \textbf{Label}\\
\midrule
1 & \texttt{elles} & \texttt{FF}\\
2 & \texttt{yamo} & \texttt{V}\\
3 & \texttt{flop} & \texttt{ADV}\\
4 & \texttt{le} & \texttt{FF}\\
5 & \texttt{kongossa} & \texttt{AMB}\\
\bottomrule
\end{tabularx}\par

\end{minipage}\hfill
\begin{minipage}[t]{0.48\linewidth}
\textbf{S06}\hfill 6 tokens

{\raggedright\texttt{go au piol shoua mes kako}\par}

\par\noindent
\begin{tabularx}{\linewidth}{@{}r X l@{}}
\toprule\tablehead
\textbf{No.} & \textbf{Raw token} & \textbf{Label}\\
\midrule
1 & \texttt{go} & \texttt{V}\\
2 & \texttt{au} & \texttt{FF}\\
3 & \texttt{piol} & \texttt{N}\\
4 & \texttt{shoua} & \texttt{V}\\
5 & \texttt{mes} & \texttt{FF}\\
6 & \texttt{kako} & \texttt{N}\\
\bottomrule
\end{tabularx}\par

\end{minipage}\par\bigskip
\begin{minipage}[t]{0.48\linewidth}
\textbf{S07}\hfill 7 tokens

{\raggedright\texttt{pere ca me chop je suis kass}\par}

\par\noindent
\begin{tabularx}{\linewidth}{@{}r X l@{}}
\toprule\tablehead
\textbf{No.} & \textbf{Raw token} & \textbf{Label}\\
\midrule
1 & \texttt{pere} & \texttt{N}\\
2 & \texttt{ca} & \texttt{FF}\\
3 & \texttt{me} & \texttt{AMB}\\
4 & \texttt{chop} & \texttt{PID}\\
5 & \texttt{je} & \texttt{FF}\\
6 & \texttt{suis} & \texttt{V}\\
7 & \texttt{kass} & \texttt{ADJ}\\
\bottomrule
\end{tabularx}\par

\end{minipage}\hfill
\begin{minipage}[t]{0.48\linewidth}
\textbf{S08}\hfill 8 tokens

{\raggedright\texttt{Ont a cote la light depuis le shap}\par}

\par\noindent
\begin{tabularx}{\linewidth}{@{}r X l@{}}
\toprule\tablehead
\textbf{No.} & \textbf{Raw token} & \textbf{Label}\\
\midrule
1 & \texttt{Ont} & \texttt{V}\\
2 & \texttt{a} & \texttt{EF}\\
3 & \texttt{cote} & \texttt{V}\\
4 & \texttt{la} & \texttt{FF}\\
5 & \texttt{light} & \texttt{N}\\
6 & \texttt{depuis} & \texttt{FF}\\
7 & \texttt{le} & \texttt{FF}\\
8 & \texttt{shap} & \texttt{ADV}\\
\bottomrule
\end{tabularx}\par

\end{minipage}\par\bigskip
}

\subsection*{Interpretation}
We obtain ADV for \code{flop} and AMB for \code{kongossa} in S05.
The latter retains competing reference labels, although we read it as a
nominal complement here. S06 places two instructions together. In S07,
we obtain N for \code{pere}, FF for \code{ca} and PID for \code{chop}.

In S08, our linguistic reading of \code{a} is a French auxiliary. The
software's English-word lookup still returns EF. We retain its V--EF--V
sequence in the computed grammar evidence and identify the contextual
auxiliary reading in this discussion.

These four statements contribute 26 occurrences. Along with the previous 25,
they account for 51 of the 81 tokens.

\clearpage
\section{Token Tables: S09--S12}
We include both rejected originals in full. We complete lexical analysis
before the parser stops at its first unsupported lookahead.

{\small\begin{minipage}[t]{0.48\linewidth}
\textbf{S09}\hfill 10 tokens

{\raggedright\texttt{shiba le price pere je n'ais pas flop les do}\par}

\par\noindent
\begin{tabularx}{\linewidth}{@{}r X l@{}}
\toprule\tablehead
\textbf{No.} & \textbf{Raw token} & \textbf{Label}\\
\midrule
1 & \texttt{shiba} & \texttt{V}\\
2 & \texttt{le} & \texttt{FF}\\
3 & \texttt{price} & \texttt{N}\\
4 & \texttt{pere} & \texttt{N}\\
5 & \texttt{je} & \texttt{FF}\\
6 & \texttt{n'ais} & \texttt{UNK}\\
7 & \texttt{pas} & \texttt{FF}\\
8 & \texttt{flop} & \texttt{ADV}\\
9 & \texttt{les} & \texttt{FF}\\
10 & \texttt{do} & \texttt{EF}\\
\bottomrule
\end{tabularx}\par

\end{minipage}\hfill
\begin{minipage}[t]{0.48\linewidth}
\textbf{S10}\hfill 9 tokens

{\raggedright\texttt{mon combi came au school alli day a pied}\par}

\par\noindent
\begin{tabularx}{\linewidth}{@{}r X l@{}}
\toprule\tablehead
\textbf{No.} & \textbf{Raw token} & \textbf{Label}\\
\midrule
1 & \texttt{mon} & \texttt{FF}\\
2 & \texttt{combi} & \texttt{N}\\
3 & \texttt{came} & \texttt{V}\\
4 & \texttt{au} & \texttt{FF}\\
5 & \texttt{school} & \texttt{N}\\
6 & \texttt{alli} & \texttt{UNK}\\
7 & \texttt{day} & \texttt{N}\\
8 & \texttt{a} & \texttt{EF}\\
9 & \texttt{pied} & \texttt{N}\\
\bottomrule
\end{tabularx}\par

\end{minipage}\par\bigskip
\begin{minipage}[t]{0.48\linewidth}
\textbf{S11}\hfill 6 tokens

{\raggedright\texttt{je falla le taco depuis bayar}\par}

\par\noindent
\begin{tabularx}{\linewidth}{@{}r X l@{}}
\toprule\tablehead
\textbf{No.} & \textbf{Raw token} & \textbf{Label}\\
\midrule
1 & \texttt{je} & \texttt{FF}\\
2 & \texttt{falla} & \texttt{V}\\
3 & \texttt{le} & \texttt{FF}\\
4 & \texttt{taco} & \texttt{N}\\
5 & \texttt{depuis} & \texttt{FF}\\
6 & \texttt{bayar} & \texttt{ADV}\\
\bottomrule
\end{tabularx}\par

\end{minipage}\hfill
\begin{minipage}[t]{0.48\linewidth}
\textbf{S12}\hfill 5 tokens

{\raggedright\texttt{les talk free le francais}\par}

\par\noindent
\begin{tabularx}{\linewidth}{@{}r X l@{}}
\toprule\tablehead
\textbf{No.} & \textbf{Raw token} & \textbf{Label}\\
\midrule
1 & \texttt{les} & \texttt{FF}\\
2 & \texttt{talk} & \texttt{V}\\
3 & \texttt{free} & \texttt{AMB}\\
4 & \texttt{le} & \texttt{FF}\\
5 & \texttt{francais} & \texttt{N}\\
\bottomrule
\end{tabularx}\par

\end{minipage}\par\bigskip
}

\subsection*{The two unresolved forms}
\code{n'ais} and \code{alli} each occur once and are the sixth token in their
respective statement. They match no explicit reference, approved annotation
or permitted morphological fallback. Each is therefore represented as UNK.

We read S09's \code{do} as money and S10's \code{a pied} as a walking
expression with a French prepositional role. Both tokens retain EF in the
unchanged program. These later labels are visible because we keep the
complete lexical output even when parsing stops after five tokens.

The lengths here are ten, nine, six and five, contributing the final 30
occurrences.

\clearpage
\section{Frequency, Coverage and Observed Variation}
We count each of the twelve statements once. We case-fold forms for frequency
counting and preserve their spelling in the token tables.
Our sample contains \CorpusTokens{} occurrences, \CorpusForms{} case-folded
forms and a mean of 6.75 tokens per statement. Lengths range from five to ten.

\begin{minipage}[t]{0.55\linewidth}
{\footnotesize\par\noindent
\begin{tabularx}{\linewidth}{@{}X l r r@{}}
\toprule
\textbf{Category} & \textbf{Code} & \textbf{Count} & \textbf{Share}\\
\midrule
FRENCH\_FUNCTION\_WORD & \texttt{FF} & 28 & 34.6\%\\
VERB & \texttt{V} & 19 & 23.5\%\\
NOUN & \texttt{N} & 14 & 17.3\%\\
ADVERB & \texttt{ADV} & 6 & 7.4\%\\
AMBIGUOUS & \texttt{AMB} & 5 & 6.2\%\\
ENGLISH\_FUNCTION\_WORD & \texttt{EF} & 3 & 3.7\%\\
ADJECTIVE & \texttt{ADJ} & 2 & 2.5\%\\
PIDGIN\_MARKER & \texttt{PID} & 2 & 2.5\%\\
UNKNOWN & \texttt{UNK} & 2 & 2.5\%\\
\bottomrule
\end{tabularx}\par
}
\par\medskip
{\small We obtain 47/81 occurrences in FF and V combined. The broad category
definitions shape this distribution. We round each percentage independently.}
\end{minipage}\hfill
\begin{minipage}[t]{0.40\linewidth}
{\small\par\noindent
\begin{tabularx}{\linewidth}{@{}X r@{}}
\toprule\tablehead
\textbf{Case-folded form} & \textbf{Occurrences}\\
\midrule
\texttt{je} & 7\\
\texttt{le} & 6\\
\texttt{depuis} & 3\\
\texttt{suis} & 3\\
\texttt{j'ai} & 2\\
\texttt{pas} & 2\\
\texttt{bayar} & 2\\
\texttt{kass} & 2\\
\texttt{go} & 2\\
\texttt{school} & 2\\
\texttt{flop} & 2\\
\texttt{au} & 2\\
\texttt{pere} & 2\\
\texttt{a} & 2\\
\texttt{les} & 2\\
\bottomrule
\end{tabularx}\par
}
\end{minipage}

\subsection*{Forty forms occur once}
{\small\texttt{n'est}, \texttt{tchop}, \texttt{nang}, \texttt{back}, \texttt{du}, \texttt{sharp}, \texttt{man}, \texttt{mban}, \texttt{gars}, \texttt{yabat}, \texttt{elles}, \texttt{yamo}, \texttt{kongossa}, \texttt{piol}, \texttt{shoua}, \texttt{mes}, \texttt{kako}, \texttt{ca}, \texttt{me}, \texttt{chop}, \texttt{ont}, \texttt{cote}, \texttt{la}, \texttt{light}, \texttt{shap}, \texttt{shiba}, \texttt{price}, \texttt{n'ais}, \texttt{do}, \texttt{mon}, \texttt{combi}, \texttt{came}, \texttt{alli}, \texttt{day}, \texttt{pied}, \texttt{falla}, \texttt{taco}, \texttt{talk}, \texttt{free}, \texttt{francais}.
}

\subsection*{Lookup coverage}
We recognize 79 of 81 occurrences (97.5\%), including ambiguous and
context-insensitive labels. This is lookup coverage rather than tagging
accuracy; we would need independent linguistic annotations to calculate
precision and recall.

We use case-folding, Unicode decomposition,
accent removal and curly-apostrophe normalization to propose groups. In this
sample we obtain only \code{J'ai} and \code{j'ai}, one occurrence each.
We discuss \code{sharp}/\code{shap} separately because their normalized
strings differ.

\clearpage
\section{Code-Mixing, Phrases and Linguistic Complexity}
We study these expressions within Cameroon's multilingual setting, where
French, English, Pidgin and indigenous languages interact \cite{ubanako}.
Ewondo and Fulfulde belong to that wider repertoire; our twelve transcriptions
do not establish the individual speakers' languages.

We obtain \CorpusSwitches{} candidate transition clues from the program's
FR, EN and PID word evidence. These clues describe lexical endpoints and
leave meaning, intonation and conversational context to our interpretation.

{\small\par\noindent
\begin{tabularx}{\linewidth}{@{}L{0.7cm} >{\raggedright\arraybackslash}X@{}}
\toprule\tablehead
\textbf{ID} & \textbf{Exact emitted transition endpoints}\\
\midrule
S01 & \texttt{pas ... tchop}; \texttt{tchop ... depuis}\\
S02 & \texttt{je ... go}\\
S03 & \texttt{suis ... back}; \texttt{back ... du}; \texttt{du ... school}; \texttt{school ... le}; \texttt{le ... sharp}; \texttt{sharp ... man}\\
S04 & None emitted\\
S05 & \texttt{le ... kongossa}\\
S06 & \texttt{go ... au}\\
S07 & \texttt{ca ... chop}; \texttt{chop ... je}\\
S08 & \texttt{la ... light}; \texttt{light ... depuis}; \texttt{le ... shap}\\
S09 & \texttt{le ... price}; \texttt{price ... pere}; \texttt{les ... do}\\
S10 & \texttt{mon ... combi}; \texttt{combi ... came}; \texttt{came ... au}; \texttt{au ... school}; \texttt{day ... pied}\\
S11 & None emitted\\
S12 & \texttt{les ... talk}; \texttt{free ... le}\\
\bottomrule
\end{tabularx}\par
}

\textbf{French frames and inserted words.}
In \code{je suis back du school}, we see a French pronoun, auxiliary and
linking form around English-origin \code{back} and \code{school}. S07 also
combines French material with \code{chop}, a Pidgin-associated form.

\textbf{Meaning changes with local use.}
We read \code{flop} in S09 as ``a lot'' and \code{do} as money. Word-by-word
English readings would miss the bargaining meaning of \code{pas flop les do}.

\textbf{Spelling is variable.}
We retain \code{sharp} in S03 and \code{shap} in S08. Both glosses refer to
the morning, suggesting a related reading that still needs contextual care.
We keep both written forms, following the distinction between speech and
transcription choices discussed by Ebongue \cite{ebongue}.

\textbf{Clauses can be juxtaposed.}
In \code{je suis kass je go nang}, we read fatigue followed by an intention
to sleep. We understand the connection although the transcription contains
neither a conjunction nor punctuation between the clauses.

\textbf{Context and intonation complete the message.}
With \code{pere ca me chop}, tone and the preceding conversation would help
us distinguish distress, complaint and emphasis. We therefore use the full
gloss and surrounding words when explaining the expression. The fixed phrase
regexes match zero verb or slang phrases in this corpus, even though these
examples clearly contain colloquial constructions.

\clearpage
\section{Constructing a Corpus-Derived Grammar}
We build our CFG from recurring lexer-category structures. The original
has ten nonterminals and 37 alternatives. For compact printing, we use
Ut, Cl, St, Bt, Pt, Lp, Mn, At, Ct and Nm to abbreviate
Utterance, Clause, SubjectTail, BareSubjectTail, PredicateTail, LinkedPhrase,
ModifiedNominal, AdverbTail, ComplementTail and Nominal respectively.

{\small\renewcommand{\arraystretch}{1.3}\par\noindent
\begin{tabularx}{\linewidth}{@{}L{0.9cm} >{\raggedright\arraybackslash}X@{}}
\toprule
\textbf{LHS} & \textbf{Original alternatives}\\
\midrule
\texttt{Ut} & \texttt{Cl} \(\mid\) \texttt{Cl} \texttt{PUNC}\\
\texttt{Cl} & \texttt{FF} \texttt{St} \(\mid\) \texttt{Nm} \texttt{Bt} \(\mid\) \texttt{V} \texttt{Pt}\\
\texttt{St} & \texttt{V} \texttt{Pt} \(\mid\) \texttt{Nm} \texttt{V} \texttt{Pt} \(\mid\) \texttt{AMB} \texttt{PID} \texttt{Ct}\\
\texttt{Bt} & \texttt{V} \texttt{Pt} \(\mid\) \texttt{FF} \texttt{St}\\
\texttt{Pt} & \texttt{V} \texttt{Pt} \(\mid\) \texttt{EF} \texttt{V} \texttt{Pt} \(\mid\) \texttt{ADJ} \texttt{Ct} \(\mid\) \texttt{AMB} \texttt{Ct} \(\mid\) \texttt{ADV} \texttt{At} \(\mid\) \texttt{Nm} \texttt{Ct} \(\mid\) \texttt{FF} \texttt{Lp} \(\mid\) \texttt{PID} \texttt{Ct} \(\mid\) \(\epsilon\)\\
\texttt{Lp} & \texttt{V} \texttt{Pt} \(\mid\) \texttt{Nm} \texttt{Ct} \(\mid\) \texttt{AMB} \texttt{Ct} \(\mid\) \texttt{ADV} \texttt{At} \(\mid\) \texttt{FF} \texttt{Mn}\\
\texttt{Mn} & \texttt{Nm} \texttt{Ct} \(\mid\) \texttt{AMB} \texttt{Ct} \(\mid\) \texttt{ADV} \texttt{At}\\
\texttt{At} & \texttt{Nm} \texttt{Ct} \(\mid\) \texttt{AMB} \texttt{Ct} \(\mid\) \texttt{FF} \texttt{Lp} \(\mid\) \texttt{V} \texttt{Pt} \(\mid\) \(\epsilon\)\\
\texttt{Ct} & \texttt{FF} \texttt{Lp} \(\mid\) \texttt{V} \texttt{Pt} \(\mid\) \(\epsilon\)\\
\texttt{Nm} & \texttt{Nm} \texttt{N} \(\mid\) \texttt{N}\\
\bottomrule
\end{tabularx}\par
}

\subsection*{Evidence-to-rule rationale}
We use Cl for a function-word subject, a nominal opening or a verb-initial
clause. We use St for \code{je suis}, \code{mon combi came} and the observed
AMB--PID pattern in \code{ca me chop}. Pt permits auxiliary/serial verbs
(\code{j'ai mban}, \code{go nang}), observed V--EF--V links, and nominal,
adjectival, adverbial or ambiguous complements.

Lp and Mn represent introduced groups such as \code{au piol},
\code{depuis bayar} and \code{depuis le shap}. Ct allows juxtaposition, as in
\code{je suis kass je go nang} and \code{go au piol shoua mes kako}.
Nm represents noun clusters; its natural left recursion is intentionally
processed by the real transformation routine.

We let Ut end with one optional punctuation token and test this extension
separately. We exclude UNKNOWN and allow AMB only in explicit positions.
The calculations here use the corpus grammar recorded with our evidence,
including its actual V--EF--V sequence for S08.

\clearpage
\section{Grammar Transformations}
\subsection*{Step 1: eliminate direct left recursion}
We detect direct left recursion in the original nonterminal Nm:
\[
  \mathrm{Nm}\rightarrow\mathrm{Nm}\ \mathrm{N}\mid\mathrm{N}.
\]
We introduce \code{Nominal\_LR1}, abbreviated Nr, and obtain:
\[
  \mathrm{Nm}\rightarrow\mathrm{N}\ \mathrm{Nr},
  \qquad
  \mathrm{Nr}\rightarrow\mathrm{N}\ \mathrm{Nr}\mid\epsilon .
\]
Both forms generate one or more noun-category tokens. The original extends
a noun chain by left recursion; the transformed form consumes the first noun
and extends it through a right-recursive suffix. No raw noun is deleted and
the language of this subgrammar is unchanged.

\subsection*{Step 2: left-factor the start rule}
The original start rule has a shared prefix:
\[
  \mathrm{Ut}\rightarrow\mathrm{Cl}\mid\mathrm{Cl}\ \mathrm{PUNC}.
\]
We introduce \code{Utterance\_LF1}, abbreviated Uf:
\[
  \mathrm{Ut}\rightarrow\mathrm{Cl}\ \mathrm{Uf},
  \qquad
  \mathrm{Uf}\rightarrow\epsilon\mid\mathrm{PUNC}.
\]
The alternative choice is postponed until the clause has been processed.
End-of-input selects epsilon; PUNC selects the optional punctuation. The
empty alternative is a suffix, not permission for an empty whole utterance.

\subsection*{Our transformation record}
We record the two operations as
\nolinkurl{eliminate_direct_left_recursion} and \nolinkurl{left_factor}.
Each record contains the complete grammar before and after the operation.
This rule ordering requires no indirect-substitution step. We also check
for residual nullable-prefix recursion.

We obtain twelve nonterminals, \GrammarProductions{} alternatives, no residual
left recursion and no LL(1) conflict. Our regression checks compare the
original and transformed languages with an independent derivation oracle
up to four terminals.

\clearpage
\section{Transformed Production Catalogue}
We assign P01--P39 in the computed grammar's order and use them consistently
in the predictive table and traces. Nr is \code{Nominal\_LR1};
Uf is \code{Utterance\_LF1}.

{\small\renewcommand{\arraystretch}{1.15}\begin{minipage}[t]{0.48\linewidth}
\par\noindent
\begin{tabularx}{\linewidth}{@{}l >{\raggedright\arraybackslash}X@{}}
\toprule\tablehead
\textbf{ID} & \textbf{Transformed production}\\
\midrule
P01 & \texttt{Ut} \(\rightarrow\) \texttt{Cl} \texttt{Uf}\\
P02 & \texttt{Cl} \(\rightarrow\) \texttt{FF} \texttt{St}\\
P03 & \texttt{Cl} \(\rightarrow\) \texttt{Nm} \texttt{Bt}\\
P04 & \texttt{Cl} \(\rightarrow\) \texttt{V} \texttt{Pt}\\
P05 & \texttt{St} \(\rightarrow\) \texttt{V} \texttt{Pt}\\
P06 & \texttt{St} \(\rightarrow\) \texttt{Nm} \texttt{V} \texttt{Pt}\\
P07 & \texttt{St} \(\rightarrow\) \texttt{AMB} \texttt{PID} \texttt{Ct}\\
P08 & \texttt{Bt} \(\rightarrow\) \texttt{V} \texttt{Pt}\\
P09 & \texttt{Bt} \(\rightarrow\) \texttt{FF} \texttt{St}\\
P10 & \texttt{Pt} \(\rightarrow\) \texttt{V} \texttt{Pt}\\
P11 & \texttt{Pt} \(\rightarrow\) \texttt{EF} \texttt{V} \texttt{Pt}\\
P12 & \texttt{Pt} \(\rightarrow\) \texttt{ADJ} \texttt{Ct}\\
P13 & \texttt{Pt} \(\rightarrow\) \texttt{AMB} \texttt{Ct}\\
P14 & \texttt{Pt} \(\rightarrow\) \texttt{ADV} \texttt{At}\\
P15 & \texttt{Pt} \(\rightarrow\) \texttt{Nm} \texttt{Ct}\\
P16 & \texttt{Pt} \(\rightarrow\) \texttt{FF} \texttt{Lp}\\
P17 & \texttt{Pt} \(\rightarrow\) \texttt{PID} \texttt{Ct}\\
P18 & \texttt{Pt} \(\rightarrow\) \(\epsilon\)\\
P19 & \texttt{Lp} \(\rightarrow\) \texttt{V} \texttt{Pt}\\
P20 & \texttt{Lp} \(\rightarrow\) \texttt{Nm} \texttt{Ct}\\
\bottomrule
\end{tabularx}\par

\end{minipage}\hfill
\begin{minipage}[t]{0.48\linewidth}
\par\noindent
\begin{tabularx}{\linewidth}{@{}l >{\raggedright\arraybackslash}X@{}}
\toprule\tablehead
\textbf{ID} & \textbf{Transformed production}\\
\midrule
P21 & \texttt{Lp} \(\rightarrow\) \texttt{AMB} \texttt{Ct}\\
P22 & \texttt{Lp} \(\rightarrow\) \texttt{ADV} \texttt{At}\\
P23 & \texttt{Lp} \(\rightarrow\) \texttt{FF} \texttt{Mn}\\
P24 & \texttt{Mn} \(\rightarrow\) \texttt{Nm} \texttt{Ct}\\
P25 & \texttt{Mn} \(\rightarrow\) \texttt{AMB} \texttt{Ct}\\
P26 & \texttt{Mn} \(\rightarrow\) \texttt{ADV} \texttt{At}\\
P27 & \texttt{At} \(\rightarrow\) \texttt{Nm} \texttt{Ct}\\
P28 & \texttt{At} \(\rightarrow\) \texttt{AMB} \texttt{Ct}\\
P29 & \texttt{At} \(\rightarrow\) \texttt{FF} \texttt{Lp}\\
P30 & \texttt{At} \(\rightarrow\) \texttt{V} \texttt{Pt}\\
P31 & \texttt{At} \(\rightarrow\) \(\epsilon\)\\
P32 & \texttt{Ct} \(\rightarrow\) \texttt{FF} \texttt{Lp}\\
P33 & \texttt{Ct} \(\rightarrow\) \texttt{V} \texttt{Pt}\\
P34 & \texttt{Ct} \(\rightarrow\) \(\epsilon\)\\
P35 & \texttt{Nm} \(\rightarrow\) \texttt{N} \texttt{Nr}\\
P36 & \texttt{Nr} \(\rightarrow\) \texttt{N} \texttt{Nr}\\
P37 & \texttt{Nr} \(\rightarrow\) \(\epsilon\)\\
P38 & \texttt{Uf} \(\rightarrow\) \(\epsilon\)\\
P39 & \texttt{Uf} \(\rightarrow\) \texttt{PUNC}\\
\bottomrule
\end{tabularx}\par

\end{minipage}\par\bigskip
}

\subsection*{How to use the catalogue}
P01 starts every parse by expanding Ut into Cl and Uf. P02--P04 select a
clause opening from FF, N or V. S02 then selects P05 and P12 for its verb and
adjective complement. Later P32 and P19 admit the juxtaposed phrase, and
P10 admits its next verb. The empty suffixes P18 and P38 complete the parse.

P35--P37 describe noun clusters. P37 can be chosen only when the next token
belongs to FOLLOW(Nr), not for every unexpected token. P38 and P39 preserve
the original optional single-punctuation choice. The table on the next two
pages makes each choice explicit rather than relying on procedural
``try alternatives until one works'' logic.

We retain full names, alternatives, generated helpers and transformation
records in the companion JSON. The abbreviations keep the printed
calculations readable.

\clearpage
\section{FIRST and FOLLOW Sets}
We calculate FIRST and FOLLOW for every transformed nonterminal.
FIRST contains possible initial terminals of a derivation and epsilon when
the nonterminal is nullable. FOLLOW contains terminals, or the end marker
\code{\$}, that can immediately follow a nonterminal in a derivation from the
start symbol. Epsilon does not belong in FOLLOW.

{\small\renewcommand{\arraystretch}{1.25}\par\noindent
\begin{tabularx}{\linewidth}{@{}L{1.8cm} >{\raggedright\arraybackslash}X >{\raggedright\arraybackslash}X@{}}
\toprule
\textbf{NT} & \textbf{FIRST} & \textbf{FOLLOW}\\
\midrule
\texttt{Ut} & \texttt{FF}, \texttt{N}, \texttt{V} & \texttt{\$}\\
\texttt{Cl} & \texttt{FF}, \texttt{N}, \texttt{V} & \texttt{\$}, \texttt{PUNC}\\
\texttt{St} & \texttt{AMB}, \texttt{N}, \texttt{V} & \texttt{\$}, \texttt{PUNC}\\
\texttt{Bt} & \texttt{FF}, \texttt{V} & \texttt{\$}, \texttt{PUNC}\\
\texttt{Pt} & \texttt{ADJ}, \texttt{ADV}, \texttt{AMB}, \texttt{EF}, \texttt{FF}, \texttt{N}, \texttt{PID}, \texttt{V}, \(\epsilon\) & \texttt{\$}, \texttt{PUNC}\\
\texttt{Lp} & \texttt{ADV}, \texttt{AMB}, \texttt{FF}, \texttt{N}, \texttt{V} & \texttt{\$}, \texttt{PUNC}\\
\texttt{Mn} & \texttt{ADV}, \texttt{AMB}, \texttt{N} & \texttt{\$}, \texttt{PUNC}\\
\texttt{At} & \texttt{AMB}, \texttt{FF}, \texttt{N}, \texttt{V}, \(\epsilon\) & \texttt{\$}, \texttt{PUNC}\\
\texttt{Ct} & \texttt{FF}, \texttt{V}, \(\epsilon\) & \texttt{\$}, \texttt{PUNC}\\
\texttt{Nm} & \texttt{N} & \texttt{\$}, \texttt{FF}, \texttt{PUNC}, \texttt{V}\\
\texttt{Nr} & \texttt{N}, \(\epsilon\) & \texttt{\$}, \texttt{FF}, \texttt{PUNC}, \texttt{V}\\
\texttt{Uf} & \texttt{PUNC}, \(\epsilon\) & \texttt{\$}\\
\bottomrule
\end{tabularx}\par
}

\subsection*{Fixed-point computation}
For each production, FIRST accumulates the first symbol's terminal evidence,
continuing through nullable prefixes and adding epsilon only when the entire
right-hand side is nullable. The implementation repeats these updates until
no set grows. FOLLOW starts by adding \code{\$} to Ut. For each nonterminal
within a right-hand side, FIRST of the suffix contributes non-epsilon
terminals; a nullable or absent suffix also contributes FOLLOW of the left
side. Those updates likewise continue to a fixed point.

\subsection*{Sanity checks on this grammar}
FIRST(Ut) is exactly \{FF, N, V\}; hence empty input cannot begin a clause.
FOLLOW(Ut) contains only end-of-input. Uf is nullable and permits punctuation,
which explains PUNC in FOLLOW(Cl). Nr can begin another N or become empty,
but its FOLLOW set is \{\$, FF, PUNC, V\}, not UNKNOWN.

We generate these sets directly from the transformed productions. The full
sets let us check epsilon entries and explain the rejection at the sixth
token in S09 and S10.

\clearpage
\section{Complete LL(1) Predictive Table}
For \(A\rightarrow\alpha\), each terminal in
\(\mathrm{FIRST}(\alpha)\setminus\{\epsilon\}\) receives that production.
If \(\alpha\) is nullable, its production is also placed under every terminal
in FOLLOW(A). Multiple distinct productions in one cell are conflicts;
the implementation retains the conflict and does not choose an arbitrary
winner.

We print the complete table below. Columns are the nine grammar terminals
plus end-of-input. The 47 populated cells refer to P01--P39 in the catalogue.
``--'' means no production and causes rejection for that nonterminal/lookahead.

{\scriptsize\setlength{\tabcolsep}{2pt}\renewcommand{\arraystretch}{1.8}
\par\noindent
\begin{tabularx}{\linewidth}{@{}l >{\centering\arraybackslash}X >{\centering\arraybackslash}X >{\centering\arraybackslash}X >{\centering\arraybackslash}X >{\centering\arraybackslash}X >{\centering\arraybackslash}X >{\centering\arraybackslash}X >{\centering\arraybackslash}X >{\centering\arraybackslash}X >{\centering\arraybackslash}X@{}}
\toprule\tablehead
\textbf{NT} & \textbf{\texttt{ADJ}} & \textbf{\texttt{ADV}} & \textbf{\texttt{AMB}} & \textbf{\texttt{EF}} & \textbf{\texttt{FF}} & \textbf{\texttt{N}} & \textbf{\texttt{PID}} & \textbf{\texttt{PUNC}} & \textbf{\texttt{V}} & \textbf{\texttt{\$}}\\
\midrule
\texttt{Ut} & -- & -- & -- & -- & P01 & P01 & -- & -- & P01 & --\\
\texttt{Cl} & -- & -- & -- & -- & P02 & P03 & -- & -- & P04 & --\\
\texttt{St} & -- & -- & P07 & -- & -- & P06 & -- & -- & P05 & --\\
\texttt{Bt} & -- & -- & -- & -- & P09 & -- & -- & -- & P08 & --\\
\texttt{Pt} & P12 & P14 & P13 & P11 & P16 & P15 & P17 & P18 & P10 & P18\\
\texttt{Lp} & -- & P22 & P21 & -- & P23 & P20 & -- & -- & P19 & --\\
\texttt{Mn} & -- & P26 & P25 & -- & -- & P24 & -- & -- & -- & --\\
\texttt{At} & -- & -- & P28 & -- & P29 & P27 & -- & P31 & P30 & P31\\
\texttt{Ct} & -- & -- & -- & -- & P32 & -- & -- & P34 & P33 & P34\\
\texttt{Nm} & -- & -- & -- & -- & -- & P35 & -- & -- & -- & --\\
\texttt{Nr} & -- & -- & -- & -- & P37 & P36 & -- & P37 & P37 & P37\\
\texttt{Uf} & -- & -- & -- & -- & -- & -- & -- & P39 & -- & P38\\
\bottomrule
\end{tabularx}\par
}

\subsection*{Worked entries}
\begin{itemize}
\item \(M[\mathrm{Ut},\mathrm{FF}]=\mathrm{P01}\), since Cl can begin with FF
and Uf follows it. The same production is selected under N and V.
\item \(M[\mathrm{Pt},\mathrm{ADJ}]=\mathrm{P12}\), selecting
\(\mathrm{Pt}\rightarrow\mathrm{ADJ}\ \mathrm{Ct}\).
\item \(M[\mathrm{Pt},\$]=\mathrm{P18}\), since the Pt alternative is empty
and end-of-input is in FOLLOW(Pt).
\item \(M[\mathrm{Nr},\mathrm{N}]=\mathrm{P36}\), whereas FF, V, PUNC and
\code{\$} select its empty alternative P37.
\item Uf selects P38 at \code{\$} and P39 at PUNC. No row has a defined
UNKNOWN entry.
\end{itemize}

\subsection*{LL(1) conclusion}
We find no competing production in a cell and no residual left recursion,
so predictive parsing is enabled. Our category grammar covers a limited
subset of everyday expressions.

\clearpage
\section{Pushdown Automaton for the LL(1) Parser}
We model our bounded parser as a deterministic pushdown recognizer.
Its extra memory is the symbol stack. For our grammar, the stack alphabet
contains twelve nonterminals, nine terminals and \code{\$}. We feed it
lexer categories and append a virtual \code{\$} to mark the end.

\automaton{automaton-parser}{5.7cm}
{Finite control of the predictive pushdown automaton. An edge label is
``input, stack top / replacement''; \code{eps} consumes no input.
The stack top is at the right. \(q[a]\) abbreviates one control state per
possible buffered lookahead, not an unbounded variable.}

\textbf{Expand, match, accept.} The initial stack becomes
\code{[\$, Ut]}. With a nonterminal \(A\) on top, the unique table entry
\(M[A,a]\) replaces it by the reversed right-hand side without consuming
the buffered category. An epsilon production only pops \(A\).
If a terminal matches \(a\), it is popped and the next lookahead is loaded.
Acceptance requires \code{\$} on top and at input end; the implementation's
final trace still shows the bottom marker rather than an empty stack.

\textbf{Why this is deterministic.} Loading reads one category into finite
control; expansion and matching then use epsilon transitions. The table has
no conflicting entry, so each valid configuration selects at most one action.
There are twenty supported lexer categories plus the end marker.
The diagram uses a prefetched lookahead; the actual code instead reads
\nolinkurl{categories[position]} and advances \code{position} only on a match.
This equivalent presentation does not change the accepted category sequences.

\textbf{Concrete execution.} In S02, the first expansion changes
\code{[\$, Ut]} to \code{[\$, Uf, Cl]}. Later matching FF consumes the first
token. S09/S10 instead reach an undefined table entry with UNK after five
matches, so neither the unknown nor the following tokens is skipped.
The complete traces on the next pages show the recorded configurations.

\begin{tabularx}{\linewidth}{@{}X r@{}}
\toprule\tablehead\textbf{Implemented guard} & \textbf{Bound}\\\midrule
Input tokens & 256\\
Characters in an individual token & 1,000\\
Recorded parse steps & 2,000\\
Stack symbols & 1,024\\
\bottomrule
\end{tabularx}

Malformed tokens and non-LL(1) grammars are rejected before the stack loop.
The implementation also checks repeated configurations and the bounds above.
Trace snapshots precede each action.

\clearpage
\section{Complete Accepted Trace: S02}
We trace \code{je suis kass je go nang}, whose emitted categories are
FF, V, ADJ, FF, V, V. The table contains every recorded step; production
identifiers refer to the full catalogue. Each stack/remaining pair precedes its
action.

{\small\renewcommand{\arraystretch}{1.2}\par\noindent
\begin{tabularx}{\linewidth}{@{}r >{\raggedright\arraybackslash}X >{\raggedright\arraybackslash}X l@{}}
\toprule
\textbf{Step} & \textbf{Stack (top at right)} & \textbf{Remaining input} & \textbf{Action}\\
\midrule
1 & \texttt{\$} \texttt{Ut} & \texttt{FF} \texttt{V} \texttt{ADJ} \texttt{FF} \texttt{V} \texttt{V} \texttt{\$} & P01\\
2 & \texttt{\$} \texttt{Uf} \texttt{Cl} & \texttt{FF} \texttt{V} \texttt{ADJ} \texttt{FF} \texttt{V} \texttt{V} \texttt{\$} & P02\\
3 & \texttt{\$} \texttt{Uf} \texttt{St} \texttt{FF} & \texttt{FF} \texttt{V} \texttt{ADJ} \texttt{FF} \texttt{V} \texttt{V} \texttt{\$} & match FF\\
4 & \texttt{\$} \texttt{Uf} \texttt{St} & \texttt{V} \texttt{ADJ} \texttt{FF} \texttt{V} \texttt{V} \texttt{\$} & P05\\
5 & \texttt{\$} \texttt{Uf} \texttt{Pt} \texttt{V} & \texttt{V} \texttt{ADJ} \texttt{FF} \texttt{V} \texttt{V} \texttt{\$} & match V\\
6 & \texttt{\$} \texttt{Uf} \texttt{Pt} & \texttt{ADJ} \texttt{FF} \texttt{V} \texttt{V} \texttt{\$} & P12\\
7 & \texttt{\$} \texttt{Uf} \texttt{Ct} \texttt{ADJ} & \texttt{ADJ} \texttt{FF} \texttt{V} \texttt{V} \texttt{\$} & match ADJ\\
8 & \texttt{\$} \texttt{Uf} \texttt{Ct} & \texttt{FF} \texttt{V} \texttt{V} \texttt{\$} & P32\\
9 & \texttt{\$} \texttt{Uf} \texttt{Lp} \texttt{FF} & \texttt{FF} \texttt{V} \texttt{V} \texttt{\$} & match FF\\
10 & \texttt{\$} \texttt{Uf} \texttt{Lp} & \texttt{V} \texttt{V} \texttt{\$} & P19\\
11 & \texttt{\$} \texttt{Uf} \texttt{Pt} \texttt{V} & \texttt{V} \texttt{V} \texttt{\$} & match V\\
12 & \texttt{\$} \texttt{Uf} \texttt{Pt} & \texttt{V} \texttt{\$} & P10\\
13 & \texttt{\$} \texttt{Uf} \texttt{Pt} \texttt{V} & \texttt{V} \texttt{\$} & match V\\
14 & \texttt{\$} \texttt{Uf} \texttt{Pt} & \texttt{\$} & P18\\
15 & \texttt{\$} \texttt{Uf} & \texttt{\$} & P38\\
16 & \texttt{\$} & \texttt{\$} & accept\\
\bottomrule
\end{tabularx}\par
}

\subsection*{Explanation of the derivation}
P02 chooses the FF-led clause. After matching \code{je}, P05 provides a verb
and predicate tail. P12 selects an adjective complement for \code{kass}.
P32 consumes the next FF and P19 allows a following verb, modeling the
juxtaposition rather than insisting on punctuation or a conjunction.

P10 permits the second verb in \code{go nang}. Once all six input categories
have been consumed, P18 removes the empty predicate tail and P38 removes the
optional punctuation suffix. The final row sees only the two end markers and
accepts. The consumed count is six and the error field is null.

\subsection*{What acceptance means}
We accept this category pattern after consuming all six tokens. Vocabulary
approval also passes because every token is recognized. Our constructed
known-word controls show why we calculate the two verdicts separately.

\clearpage
\section{Complete Rejected Traces: S09 and S10}
We complete lexical analysis of both inputs, then reject at the first
unavailable parsing-table entry. We preserve the unknown token in the trace.

\textbf{S09: \code{n'ais} at position 6.}
{\scriptsize\renewcommand{\arraystretch}{1.02}\par\noindent
\begin{tabularx}{\linewidth}{@{}r >{\raggedright\arraybackslash}X >{\raggedright\arraybackslash}X l@{}}
\toprule\tablehead
\textbf{Step} & \textbf{Stack (top at right)} & \textbf{Remaining input} & \textbf{Action}\\
\midrule
1 & \texttt{\$} \texttt{Ut} & \texttt{V} \texttt{FF} \texttt{N} \texttt{N} \texttt{FF} \texttt{UNK} \texttt{FF} \texttt{ADV} \texttt{FF} \texttt{EF} \texttt{\$} & P01\\
2 & \texttt{\$} \texttt{Uf} \texttt{Cl} & \texttt{V} \texttt{FF} \texttt{N} \texttt{N} \texttt{FF} \texttt{UNK} \texttt{FF} \texttt{ADV} \texttt{FF} \texttt{EF} \texttt{\$} & P04\\
3 & \texttt{\$} \texttt{Uf} \texttt{Pt} \texttt{V} & \texttt{V} \texttt{FF} \texttt{N} \texttt{N} \texttt{FF} \texttt{UNK} \texttt{FF} \texttt{ADV} \texttt{FF} \texttt{EF} \texttt{\$} & match V\\
4 & \texttt{\$} \texttt{Uf} \texttt{Pt} & \texttt{FF} \texttt{N} \texttt{N} \texttt{FF} \texttt{UNK} \texttt{FF} \texttt{ADV} \texttt{FF} \texttt{EF} \texttt{\$} & P16\\
5 & \texttt{\$} \texttt{Uf} \texttt{Lp} \texttt{FF} & \texttt{FF} \texttt{N} \texttt{N} \texttt{FF} \texttt{UNK} \texttt{FF} \texttt{ADV} \texttt{FF} \texttt{EF} \texttt{\$} & match FF\\
6 & \texttt{\$} \texttt{Uf} \texttt{Lp} & \texttt{N} \texttt{N} \texttt{FF} \texttt{UNK} \texttt{FF} \texttt{ADV} \texttt{FF} \texttt{EF} \texttt{\$} & P20\\
7 & \texttt{\$} \texttt{Uf} \texttt{Ct} \texttt{Nm} & \texttt{N} \texttt{N} \texttt{FF} \texttt{UNK} \texttt{FF} \texttt{ADV} \texttt{FF} \texttt{EF} \texttt{\$} & P35\\
8 & \texttt{\$} \texttt{Uf} \texttt{Ct} \texttt{Nr} \texttt{N} & \texttt{N} \texttt{N} \texttt{FF} \texttt{UNK} \texttt{FF} \texttt{ADV} \texttt{FF} \texttt{EF} \texttt{\$} & match N\\
9 & \texttt{\$} \texttt{Uf} \texttt{Ct} \texttt{Nr} & \texttt{N} \texttt{FF} \texttt{UNK} \texttt{FF} \texttt{ADV} \texttt{FF} \texttt{EF} \texttt{\$} & P36\\
10 & \texttt{\$} \texttt{Uf} \texttt{Ct} \texttt{Nr} \texttt{N} & \texttt{N} \texttt{FF} \texttt{UNK} \texttt{FF} \texttt{ADV} \texttt{FF} \texttt{EF} \texttt{\$} & match N\\
11 & \texttt{\$} \texttt{Uf} \texttt{Ct} \texttt{Nr} & \texttt{FF} \texttt{UNK} \texttt{FF} \texttt{ADV} \texttt{FF} \texttt{EF} \texttt{\$} & P37\\
12 & \texttt{\$} \texttt{Uf} \texttt{Ct} & \texttt{FF} \texttt{UNK} \texttt{FF} \texttt{ADV} \texttt{FF} \texttt{EF} \texttt{\$} & P32\\
13 & \texttt{\$} \texttt{Uf} \texttt{Lp} \texttt{FF} & \texttt{FF} \texttt{UNK} \texttt{FF} \texttt{ADV} \texttt{FF} \texttt{EF} \texttt{\$} & match FF\\
14 & \texttt{\$} \texttt{Uf} \texttt{Lp} & \texttt{UNK} \texttt{FF} \texttt{ADV} \texttt{FF} \texttt{EF} \texttt{\$} & reject\\
\bottomrule
\end{tabularx}\par
}

\textbf{S10: \code{alli} at position 6.}
{\scriptsize\renewcommand{\arraystretch}{1.02}\par\noindent
\begin{tabularx}{\linewidth}{@{}r >{\raggedright\arraybackslash}X >{\raggedright\arraybackslash}X l@{}}
\toprule
\textbf{Step} & \textbf{Stack (top at right)} & \textbf{Remaining input} & \textbf{Action}\\
\midrule
1 & \texttt{\$} \texttt{Ut} & \texttt{FF} \texttt{N} \texttt{V} \texttt{FF} \texttt{N} \texttt{UNK} \texttt{N} \texttt{EF} \texttt{N} \texttt{\$} & P01\\
2 & \texttt{\$} \texttt{Uf} \texttt{Cl} & \texttt{FF} \texttt{N} \texttt{V} \texttt{FF} \texttt{N} \texttt{UNK} \texttt{N} \texttt{EF} \texttt{N} \texttt{\$} & P02\\
3 & \texttt{\$} \texttt{Uf} \texttt{St} \texttt{FF} & \texttt{FF} \texttt{N} \texttt{V} \texttt{FF} \texttt{N} \texttt{UNK} \texttt{N} \texttt{EF} \texttt{N} \texttt{\$} & match FF\\
4 & \texttt{\$} \texttt{Uf} \texttt{St} & \texttt{N} \texttt{V} \texttt{FF} \texttt{N} \texttt{UNK} \texttt{N} \texttt{EF} \texttt{N} \texttt{\$} & P06\\
5 & \texttt{\$} \texttt{Uf} \texttt{Pt} \texttt{V} \texttt{Nm} & \texttt{N} \texttt{V} \texttt{FF} \texttt{N} \texttt{UNK} \texttt{N} \texttt{EF} \texttt{N} \texttt{\$} & P35\\
6 & \texttt{\$} \texttt{Uf} \texttt{Pt} \texttt{V} \texttt{Nr} \texttt{N} & \texttt{N} \texttt{V} \texttt{FF} \texttt{N} \texttt{UNK} \texttt{N} \texttt{EF} \texttt{N} \texttt{\$} & match N\\
7 & \texttt{\$} \texttt{Uf} \texttt{Pt} \texttt{V} \texttt{Nr} & \texttt{V} \texttt{FF} \texttt{N} \texttt{UNK} \texttt{N} \texttt{EF} \texttt{N} \texttt{\$} & P37\\
8 & \texttt{\$} \texttt{Uf} \texttt{Pt} \texttt{V} & \texttt{V} \texttt{FF} \texttt{N} \texttt{UNK} \texttt{N} \texttt{EF} \texttt{N} \texttt{\$} & match V\\
9 & \texttt{\$} \texttt{Uf} \texttt{Pt} & \texttt{FF} \texttt{N} \texttt{UNK} \texttt{N} \texttt{EF} \texttt{N} \texttt{\$} & P16\\
10 & \texttt{\$} \texttt{Uf} \texttt{Lp} \texttt{FF} & \texttt{FF} \texttt{N} \texttt{UNK} \texttt{N} \texttt{EF} \texttt{N} \texttt{\$} & match FF\\
11 & \texttt{\$} \texttt{Uf} \texttt{Lp} & \texttt{N} \texttt{UNK} \texttt{N} \texttt{EF} \texttt{N} \texttt{\$} & P20\\
12 & \texttt{\$} \texttt{Uf} \texttt{Ct} \texttt{Nm} & \texttt{N} \texttt{UNK} \texttt{N} \texttt{EF} \texttt{N} \texttt{\$} & P35\\
13 & \texttt{\$} \texttt{Uf} \texttt{Ct} \texttt{Nr} \texttt{N} & \texttt{N} \texttt{UNK} \texttt{N} \texttt{EF} \texttt{N} \texttt{\$} & match N\\
14 & \texttt{\$} \texttt{Uf} \texttt{Ct} \texttt{Nr} & \texttt{UNK} \texttt{N} \texttt{EF} \texttt{N} \texttt{\$} & reject\\
\bottomrule
\end{tabularx}\par
}

S09 ends at Lp (LinkedPhrase) with lookahead UNK; its legal lookaheads are
ADV, AMB, FF, N and V. S10 ends at Nr (Nominal\_LR1) with UNK; legal
lookaheads are \code{\$}, FF, N, PUNC and V. Both have consumed five tokens.
The different stack tops explain the distinct error messages even though the
failure positions match.

\clearpage
\section{Evaluation Results and Boundary Controls}
We evaluate all twelve originals with the same grammar and reviewed lexicon
snapshot. We compute every token, verdict, consumed count and trace length
shown below.

{\small\renewcommand{\arraystretch}{1.2}\par\noindent
\begin{tabularx}{\linewidth}{@{}l r X X r r@{}}
\toprule
\textbf{ID} & \textbf{Tokens} & \textbf{Vocabulary} & \textbf{CFG} & \textbf{Consumed} & \textbf{Trace rows}\\
\midrule
S01 & 6 & Pass & Accept & 6 & 18\\
S02 & 6 & Pass & Accept & 6 & 16\\
S03 & 8 & Pass & Accept & 8 & 22\\
S04 & 5 & Pass & Accept & 5 & 14\\
S05 & 5 & Pass & Accept & 5 & 14\\
S06 & 6 & Pass & Accept & 6 & 20\\
S07 & 7 & Pass & Accept & 7 & 19\\
S08 & 8 & Pass & Accept & 8 & 21\\
S09 & 10 & Fail & Reject & 5 & 14\\
S10 & 9 & Fail & Reject & 5 & 14\\
S11 & 6 & Pass & Accept & 6 & 18\\
S12 & 5 & Pass & Accept & 5 & 16\\
\bottomrule
\end{tabularx}\par
}

We accept ten category sequences (83.3\%) and reject two (16.7\%).
We then use the following constructed controls to check the distinction
between vocabulary approval and grammar acceptance.

{\small\renewcommand{\arraystretch}{1.2}\par\noindent
\begin{tabularx}{\linewidth}{@{}l X l l@{}}
\toprule\tablehead
\textbf{ID} & \textbf{Constructed control} & \textbf{Vocabulary} & \textbf{CFG}\\
\midrule
C01 & (empty) & Pass & Reject\\
C02 & \texttt{school} & Pass & Reject\\
C03 & \texttt{je} & Pass & Reject\\
C04 & \texttt{je suis kass.} & Pass & Accept\\
C05 & \texttt{je suis kass?!} & Pass & Reject\\
C06 & \texttt{je je suis} & Pass & Reject\\
C07 & \texttt{je suis alli} & Fail & Reject\\
\bottomrule
\end{tabularx}\par
}

Empty input has no UNKNOWN but cannot derive a Clause. A noun-only fragment
and a dangling FF-led clause also reject. One punctuation token is supported;
two are not. Known categories in an unsupported order can reject without any
unknown word.

We also test unseen combinations, unsupported category orders and bounded
language equivalence. A larger independently annotated corpus would help us
assess how well the grammar generalizes.

\clearpage
\section{Working Analyzer: Browser Evidence}
We illustrate the workflow with captures from
\url{https://camfranglais.duckdns.org} dated 27 September 2026. We inspected
an existing saved test and its original results.

\screen{public-test-20260927.png}{11.8cm}
{Existing S02 selected through Analysis, with its original input, vocabulary
verdict and token details.}

\screen{public-analyzer-20260927.png}{6.5cm}
{The current analyzer opens directly without login. Submission explicitly warns
that tests are public and retained. This crop shows the input panel.}

We retrieve S02's original token table and trace when selecting its saved
test. The capture register and uncropped images accompany the report.

\clearpage
\section{Topic Patterns in the Corpus}
We relate each statement to an everyday theme using its wording and French
gloss. The table gives the exact cues behind our reading, making it clear
how we arrived at each topic. These are our interpretations of the text,
rather than an additional output of the application.

{\small\renewcommand{\arraystretch}{1.35}
\begin{tabularx}{\linewidth}{@{}L{0.7cm}L{4.4cm}X@{}}
\toprule\tablehead\textbf{ID} & \textbf{Textual cues} &
\textbf{Our topic reading}\\\midrule
S01 & \code{tchop} & Food and eating\\
S02 & \code{kass, nang} & Fatigue and rest\\
S03 & \code{back, school} & School and a return journey\\
S04 & \code{mban, yabat} & Setback and personal distress\\
S05 & \code{kongossa} & Gossip and social talk\\
S06 & \code{piol, shoua, kako} & Household tasks\\
S07 & \code{kass} & Fatigue and personal state\\
S08 & \code{cote, light} & Electricity interruption\\
S09 & \code{shiba, price, do} & Market bargaining and transactions\\
S10 & \code{came, school, a pied} & Walking to school; commuting\\
S11 & \code{falla, taco} & Seeking transport\\
S12 & \code{talk, francais} & Language use\\
\bottomrule
\end{tabularx}}

\subsection*{How we read the patterns}
We pair \code{cote} with \code{light} for electricity and
\code{shiba}, \code{price} and \code{do} for bargaining.
\code{school} suggests university or school life; movement words and
\code{a pied} add the commuting reading. We read \code{taco} as transport
because the corresponding gloss explicitly mentions a taxi.

We use surrounding words to handle ambiguous cues. In S02, \code{go} occurs
with sleeping and fatigue, so we prefer rest over a transport label.
Security and weather remain unrepresented in these twelve statements.
Our analysis covers the observed themes without adding new raw statements.

\subsection*{Reference material}
We consult Collection for the original entries and Dictionary for
source-labelled meanings. We keep the 26 synthetic practice examples
separate from this corpus. The two screenshots in Section~22 illustrate
the implemented analyzer, while the table above presents our topic readings.

\clearpage
\section{Application Architecture}
We organize the application into four layers. The interface receives a
statement, the backend coordinates the request, the compiler produces the
lexical and grammar results, and storage preserves the test. We describe
the detailed classes and interactions in the separate SDD and UML atlas.

\begin{tabularx}{\linewidth}{@{}L{3.0cm}>{\raggedright\arraybackslash}X@{}}
\toprule\tablehead\textbf{Layer} & \textbf{Responsibility}\\\midrule
Frontend & React/TypeScript pages for input, results, Collection, Dictionary
and synthetic examples.\\
Backend & FastAPI request handling, analysis coordination and retrieval of
saved results.\\
Lexer / parser & Regex and vocabulary rules, frequencies, grammar
transformations, FIRST/FOLLOW, the LL(1) table and parsing traces.\\
Storage & SQLite records for Collection, the saved grammar and original
test snapshots.\\
\bottomrule
\end{tabularx}

\section{Verification, Reproducibility and Limitations}
\subsection*{Our checks}
We compare the report with a fresh computation from the unchanged twelve
statements. This revision improves the writing and linguistic discussion
while preserving the application, grammar and measured compiler results.

\begin{tabularx}{\linewidth}{@{}L{3.1cm}X@{}}
\toprule\tablehead\textbf{Evidence} & \textbf{Observed result}\\\midrule
Raw corpus & All twelve strings match their fixtures exactly.\\
Lexical output & All 81 tokens, frequencies, category counts and variations
match the recorded evidence.\\
Grammar & The two transformations yield 39 productions, 47 populated
LL(1) cells and zero conflicts.\\
Parser & Ten originals accept; S09 and S10 reject at token six.
The complete traces and constructed controls match their expected results.\\
Report & Generated tables, embedded figures, references and PDF layout
are checked against the source material.\\
\bottomrule
\end{tabularx}

\subsection*{Reproduction procedure}
We run \code{python -m tools.build\_report\_evidence --check}, the corpus
regression tests, the local LaTeX build and
\code{python -m tools.check\_documentation}. The JSON retains full grammar
names, intermediate transformations and traces.

Our main limits are the small development corpus, broad lexical categories
and unresolved forms. The linguistic readings in Sections~7 and 23 help us
explain those limits without replacing the measured outputs.

\clearpage
\section{Conclusion}
We have built and examined a complete compiler pipeline around twelve
unchanged transcriptions. We classify all 81 token occurrences, compute
frequencies and language clues, transform the grammar and expose the
conflict-free LL(1) table. The traces show how each lexical category affects
the next parsing action.

We accept ten statements and locate the unresolved forms in S09 and S10
at token six. Our controls also show that recognized words can occur in
an unsupported category order. Through the slang and topic readings,
we relate the formal output to everyday meaning, variable spelling
and juxtaposed clauses.

We present these stages in a public interface that keeps original test
results available for inspection. The project has helped us connect regex
recognition, grammar transformations and predictive parsing with the
complexity of informal communication.

\subsection*{Submission and presentation}
We identify \code{tchop}, the contextual uses of \code{a}, and the two unknown
forms as priorities for later lexical review. We preserve the current
compiler results so that the report's tables and traces remain reproducible.

Our deadline is \textbf{29 September 2026 at 14:00}. In the examination,
\textbf{each student presents for three minutes}. We will rehearse our
individual parts and use an existing test to explain the main workflow.

\subsection*{Author group}
We present this work jointly as the three-member group below. We share
credit for entering the twelve statements into the compiler and preparing
the corpus used in this report.

{\small\renewcommand{\arraystretch}{1.6}
\begin{tabularx}{\linewidth}{@{}X r@{}}
\toprule\tablehead\textbf{Member} & \textbf{Matricule}\\\midrule
Kwete Ngouba Junior Rayan & ICTU20241377\\
Djemtchimo Noukui Bruno Jonatan & ICTU20241585\\
Amina Boubakary & ICTU20241870\\
\bottomrule
\end{tabularx}}

\clearpage
\addcontentsline{toc}{section}{References and Evidence Register}
\renewcommand{\refname}{References and Evidence Register}
\begin{thebibliography}{9}
\bibitem{brief}
Faculty of Information and Communication Technologies.
\emph{CS4110 Compiler Construction, Summer 2026, SET A:
Lexical and Syntactic Analysis of Informal Urban Communication in Yaound\'e}.
Assignment brief; instructor Engr. Tanwi Nkiamboh; assessment dated
29 September 2026. Source of the 10--15-statement, group-of-three and
25--30-page requirements.

\bibitem{ebongue}
Augustin Emmanuel Ebongue.
``Linguistic Features in a Marginal Corpus. The Case of Written Camfranglais,
a Cameroon Youth Language.''
\emph{Multilingual Margins}, 9(1), issue year 2022; published online
25 May 2023. Publisher metadata and abstract checked 27 September 2026.
\url{https://www.epubs.ac.za/index.php/mm/article/view/1392}.

\bibitem{ubanako}
Valentine Njende Ubanako.
``Is `Camfranglais' A New Language? A Review of Current Opinions.''
\emph{International Linguistics Research}, 4(1), 2021.
We use its discussion of Cameroon's wider language-contact setting.
\url{https://doi.org/10.30560/ilr.v4n1p36}.

\bibitem{compiler}
Alfred V. Aho, Monica S. Lam, Ravi Sethi and Jeffrey D. Ullman.
\emph{Compilers: Principles, Techniques, and Tools}, second edition.
Pearson/Addison-Wesley. General compiler background for lexical analysis,
grammar transformations and predictive parsing.

\bibitem{source}
Camfranglais project source, baseline \code{a1b88aa}.
\nolinkurl{compiler/lexer/tokenizer.py},
\nolinkurl{compiler/lexer/languages.py},
\nolinkurl{compiler/lexer/reference.py},
\nolinkurl{compiler/parser/yaounde.py},
\nolinkurl{compiler/parser/analysis.py},
\nolinkurl{compiler/parser/predictive.py}.
We reproduce the reported outputs from these implemented rules.

\bibitem{corpus}
Project Collection, twelve exact transcriptions S01--S12.
\nolinkurl{docs/evidence/final-corpus-20260927.json}.
Read-only analysis snapshot dated 27 September 2026, preserving raw wording
and French glosses. The initial compiler-input record is dated
22 September 2026; Section~2 describes our corpus preparation.

\bibitem{computed}
Reproducible coursework analysis.
\nolinkurl{tools/build_report_evidence.py},
\nolinkurl{docs/evidence/final-analysis-20260927.json},
\nolinkurl{docs/report-data/}.
Contains source fingerprints, full tokens, sets, table, transformations,
traces and boundary controls. Tested against the twelve exact cases in
\nolinkurl{compiler/tests/yaounde_cases.py}.

\bibitem{reference}
Classified vocabulary and additional reference senses.
\nolinkurl{dictionary/full_lexicon_classified.csv},
\nolinkurl{dictionary/camfranglais.md},
\nolinkurl{dictionary/extra_lexicon.md}.
Source-labelled lookup material: \ReferenceRows{} CSV rows and
\DictionaryCount{} dictionary display entries.

\bibitem{screens}
Public application screenshots, captured 27 September 2026.
\nolinkurl{docs/evidence/screenshots/};
\url{https://camfranglais.duckdns.org}.
We inspected existing records without posting new tests. The capture
register records image hashes, context and the saved-test crop.
\end{thebibliography}

Supplementary SRS, SDD and UML atlas sources and PDFs are distributed in the
same documentation directory.
\end{document}
